\documentclass[11pt]{article}
\usepackage{fontspec}
\usepackage[french,english]{babel}
\usepackage{amsmath}
\usepackage{amssymb}
\usepackage{amsthm}
\usepackage{mathrsfs}
\usepackage{tikz-cd}
\usepackage{enumitem}
\usetikzlibrary{babel}

\newcommand{\printedpage}[1]{\marginpar{\footnotesize p.~#1}}
\newcommand{\msp}[1]{\marginpar{\footnotesize [#1]}}
\newcommand{\HOT}{\mathrm{HOT}}
\newcommand{\hot}{\mathrm{hot}}
\newcommand{\Cat}{\mathrm{Cat}}
\newcommand{\id}{\mathrm{id}}
\newcommand{\pr}{\mathrm{pr}}
\newcommand{\Ob}{\mathrm{Ob}}
\newcommand{\Wd}{W^{\mathrm{d}}}

\setlength{\parindent}{0pt}
\setlength{\parskip}{0.6em}

\begin{document}

\printedpage{84}
% head: A. Grothendieck : Dérivateurs, Ch. VI, Édition des Universités de Montpellier II et Paris VII
even better still for the functor $f_!$ between the HOT, since the description of this functor can be given directly on the whole of $\Cat/S'$, without having to restrict oneself to more or less \msp{page 138}clever subcategories, where the induced functor would be compatible with the localizers ($\underline{\mathrm{Liss}}\,S$ and $\underline{\mathrm{Fibhot}}\,S$ in the case of $f^*$).

We have immediately the transitivity of $f_!$, $f_!^{\mathrm{lc}}$, $f_!^{o}$, i.e.
\[
(gf)_! = g_!f_!, \qquad (gf)_!^{\mathrm{lc}} = g_!^{\mathrm{lc}}f_!^{\mathrm{lc}}, \qquad (gf)_!^{o} = g_!^{o}f_!^{o}.
\]

But we are not going to move $S$, $S'$, $f$ for the moment, for we must now, before anything else, try to establish the adjunction of $f_!$ and $f^*$ (in that order).

For the moment, let us deal with $\HOT_W$ only - the case of $\HOT_W^{o}$ will follow from this by applying what we shall know to $f^{o}$. And we shall deal with $\HOT_W^{\mathrm{lc}}$ later.

1) \textit{The homomorphism}
\[
\alpha : \id_{\HOT_W(S')} \longrightarrow f^*f_!.
\]

If $X'$ is in $\Cat/S'$, $f_!(\hot_{S',W}(X'))$ is $\hot_{S,W}(X')$, but we cannot find its inverse image $f^*$ over $S'$ by taking the categorical inverse image. For even if we make draconian hypotheses on $X'$ as an object of $\Cat/S'$ (e.g. $X'$ fibered over $S'$), these will be reduced to nothing when we have composed with $f$, on which we do not want to make any hypothesis here. We \msp{page 139}must therefore first apply the functor $\Phi_S$ to $X'$ over $S$, then take its inverse image over $S'$:
\[
f^*f_!(\hot_{S',W}(X')) = \hot_{S',W}(\Phi_S(X'/S)\times_S S').
\]
We must therefore find an arrow
\[
\alpha_{X'} : X' \longrightarrow \Phi_S(X'/S)\times_S S'
\]
in the localized $(\Cat/S')(\Wd_{S'})^{-1}$, or even an arrow in $\Cat/S'$ itself. But by the adjunction properties of the functors $f_!$ and $f^*$ at the \textit{categorical level,} such an arrow amounts to an arrow
\[
f_!(X') \longrightarrow \Phi_S(X'/S)
\]
in $\Cat/S$, i.e. to an arrow which makes the square commute
\[
\begin{tikzcd}
X' \arrow[r] \arrow[d] & \Phi_S(X'/S') \arrow[d] \\ % [sic] \Phi_S(X'/S') -> \Phi_S(X'/S), the arrow f_!(X') -> \Phi_S(X'/S) just above has X'/S
S' \arrow[r] & S.
\end{tikzcd}
\]
% foot: Version 9 novembre 2001

\printedpage{85}
% head: A. Grothendieck : Dérivateurs, Ch. VI, Édition des Universités de Montpellier II et Paris VII
Now we have such a canonical arrow, namely $i_{X'/S}$. It is functorial in $X'$ over $S$ (without even having to worry whether $X'$ over $S$ comes from an $X'$ over $S'$, i.e. whether the structural arrow $X'\longrightarrow S$ factors through $S'\longrightarrow S$).

\leavevmode\msp{page 140}2) \textit{The homomorphism}
\[
\beta : \id_{\HOT_W(S)} \longleftarrow f_!f^*.
\]

We can describe $\HOT_W(S)$ by $(\underline{\mathrm{Fib}}\,S)(\Wd_S)^{-1}$ or $(\underline{\mathrm{Liss}}\,S)(\Wd_S)^{-1}$ for example, so that $f^*$ has a simple description by
\[
f^*(\hot_{S,W}(X)) = \hot_{S',W}(X\times_S S')
\]
(for $X$ smooth over $S$). We then have
\[
f_!f^*(\hot_{S,W}(X)) = \hot_{S,W}(X\times_S S'),
\]
but watch out, $X\times_S S'$ is not a priori in $\underline{\mathrm{Liss}}\,S$, only in $\Cat/S$ and nothing more. But that does not prevent it from giving what is needed, by localization $(\Wd_S)^{-1}$ in $\Cat/S$. To define $\beta_X$, it therefore suffices to find an arrow
\[
\beta_X : X \longleftarrow X\times_S S'
\]
in $\Cat/S$, and we do see $\pr_1^X$ (\footnote[68]{when $\xi\in\HOT_W(S)$ is defined by $X$ in $\Cat/S$ not necessarily smooth, we must find $\Phi_S(X)\times_S S'\longrightarrow X$ in $(\Cat/S)(\Wd_S)^{-1}$, we take $(i_X)^{-1}\pr_1^{\Phi_S(X)}$ (cf. page 143).}). We thus have
\[
\left\{
\begin{array}{lcl}
\alpha : \id & \longrightarrow & f^*f_! \\
\beta : \id & \longleftarrow & f_!f^*
\end{array}
\right.
\]

3) \textit{The compatibility}
\[
\begin{tikzcd}
f_! \arrow[r, "f_!*\alpha"] \arrow[rr, bend right, "\id"'] & f_!f^*f_! \arrow[r, "\beta*f_!"] & f_!.
\end{tikzcd}
\]

We must therefore take $X'$ over $S'$, and make the arrows explicit starting from the object
\[
\xi' = \hot_{S',W}(X').
\]
% foot: Version 9 novembre 2001

\printedpage{86}
% head: A. Grothendieck : Dérivateurs, Ch. VI, Édition des Universités de Montpellier II et Paris VII
\leavevmode\msp{page 141}We have
\[
f_!(\xi') = \hot_{S,W}(X'/S),
\]
and to calculate $f^*f_!(\xi')$, we must as before take $\Phi_S(X')$ and we find
\[
f^*f_!(\xi') = \hot_{S',W}(\Phi_S(X')\times_S S').
\]
Finally,
\[
f_!f^*f_!(\xi') = \hot_{S,W}(\Phi_S(X')\times_S S').
\]
We must make the arrows explicit
\[
\begin{tikzcd}[row sep=small]
\underbrace{f_!(\xi')} \arrow[r, "{f_!(\alpha_{\xi'})}"] & \underbrace{f_!f^*f_!(\xi')} \arrow[r, "{\beta_{f_!(\xi')}}"] & \underbrace{f_!(\xi')}. \\
X'/S \arrow[r, "{\alpha_{X'}}", "{\text{deduced from } i_{X'}}"'] & \Phi_S(X')\times_S S' \arrow[r, dotted, "{\beta_{X'}}"] & X'/S \\
 & \text{over } S &
\end{tikzcd}
\tag*{$(*)$}
\]

$\alpha_{X'}$ is deduced from $i_{X'} : X'\longrightarrow\Phi_S(X')$ as the unique $S'$-morphism making the diagram commute
\[
\begin{tikzcd}
X' \arrow[r, "{\alpha_{X'}}"] \arrow[d] \arrow[rr, bend left, "{i_{X'}}"] & \Phi_S(X')\times_S S' \arrow[r, "\pr_1"] \arrow[dl, dotted] & \Phi_S(X') \arrow[d] \\ % Reading verified against PDF leaf 48: the curved i_X' arrow is solid.
S' \arrow[rr] & & S.
\end{tikzcd}
\tag*{$(**)$}
\]

To describe $\beta_{f_!(\xi')}$, where $f_!(\xi') = \hot_{S,W}(X')$, we must give a general description of $\beta_\xi$ that is valid for \textit{every} representative $X$ of $\xi$, not only in the case $X$ smooth over $S$. To calculate $f^*(\xi)$, we must therefore take the $\hot_{S',W}$ of $\Phi(X)\times_S S'$, then for $f_!f^*(\xi)$, we \msp{page 142}must take the $\hot_{S,W}$ of the same $\Phi_S(X)\times_S S'$. We must define, if we can, an arrow
\[
\Phi_S(X)\times_S S' \longrightarrow X \qquad\qquad \text{in } \Cat/S,
\]
or at least in $(\Cat/S)(\Wd_S)^{-1}$. But we have
\[
i_X : X \longrightarrow \Phi_S(X)
\]
which is in $\Wd_S$, and which by Cartesian product defines
\[
i_X\times_S\id_{S'} : X\times_S S' \longrightarrow \Phi_S(X)\times_S S'
\]
(in other words we have taken the inverse image of $i_X$ under $S'\longrightarrow S$). Is this arrow still in $\Wd_S$? (\footnote[69]{No, it is not, for if $S'=\{s\}$, this would mean that $X'_s\longrightarrow s\backslash X$ is in $W$, which is false in general.}). If so, we have the diagram
% foot: Version 9 novembre 2001
\printedpage{87}
% head: A. Grothendieck : Dérivateurs, Ch. VI, Édition des Universités de Montpellier II et Paris VII
\[
\begin{tikzcd}
\Phi_S(X)\times_S S' \arrow[r, dotted, "\beta_X ?"] & X \\
X\times_S S' \arrow[u, "{i_X\times_S\id_{S'}}", "{=\pi_X}"'] \arrow[ur, "\pr_1^X"] &
\end{tikzcd}
\]
and we define the desired arrow $\beta_X$ \textit{in} $(\Cat/S)(\Wd_S)^{-1}$ by the commutativity of this triangle, thus as
\[
\beta_X = \pr_1^X\circ\underbrace{\pi_X^{-1}}_{\text{inverse in }(\Cat/S)(\Wd_S)^{-1}} \qquad\qquad .
\]

\leavevmode\msp{page 143}But that doesn't work - cf. annotation! We shall then do the reverse
\[
\begin{tikzcd}
\Phi_S(X)\times_S S' \arrow[r, "\pr_1^{\Phi_S(X)}"] \arrow[dr, dotted, "\beta_X ?"'] & \Phi_S(X) \\
 & X \arrow[u, "{i_X \qquad \text{in } \Wd_S}"']
\end{tikzcd}
\]

(\footnote[70]{Description of $\beta_\xi$ for $\xi=\hot_{S,W}(X)$, with $X$ arbitrary in $\Cat/S$, to be made explicit in definition in due form.}), we define $\beta_X$ as $(i_X)^{-1}\pr_1^{\Phi_S(X)}$, where $(i_X)^{-1}$ is defined in $(\Cat/S)(\Wd_S)^{-1}$.

To prove the commutativity of $(*)$ (p. 141), since $i_{X'} : X'\longrightarrow\Phi_S(X')$ is in $\Wd_S$, it suffices to prove, instead of $\beta_{X'}\alpha_{X'} = \id_{X'}$, the relation deduced by composing with $i_{X'} : X'\longrightarrow\Phi_S(X')$
\[
\underbrace{i_{X'}\beta_{X'}}_{\pr_1^{\Phi_S(X')}}\alpha_{X'} = i_{X'},
\]
i.e.
\[
\pr_1^{\Phi_S(X')}\circ\alpha_{X'} = i_{X'},
\]
which is nothing other than the defining relation of $\alpha_{X'}$ (cf. $(**)$ p. 141), we win!

4) \textit{Compatibility}
\[
\begin{tikzcd}
f^* \arrow[r, "\alpha*f^*"] \arrow[rr, bend right, "\id"'] & f^*f_!f^* \arrow[r, "f^**\beta"] & f^*.
\end{tikzcd}
\]
% foot: Version 9 novembre 2001

\printedpage{88}
% head: A. Grothendieck : Dérivateurs, Ch. VI, Édition des Universités de Montpellier II et Paris VII
\leavevmode\msp{page 144}Let us therefore start with a $\xi\in\Ob\HOT_W(S)$, which we may suppose defined by $X\in\Ob\underline{\mathrm{Fib}}\,S$. Then
\[
\begin{aligned}
f^*(\xi) &\;=\; \hot_{S',W}(X\times_S S'),\\
f_!f^*(\xi) &\;=\; \hot_{S,W}(X\times_S S'),\\
f^*f_!f^*(\xi) &\;=\; \hot_{S',W}(\Phi_S(X\times_S S')\times_S S').
\end{aligned}
\]
We must now make explicit
\[
\underbrace{f^*(\xi)}_{\hot_{S',W}(X')} \xrightarrow{\alpha_{f^*(\xi)}} \underbrace{f^*f_!f^*(\xi)}_{\hot_{S',W}(\Phi_S(X')\times_S S')} \xrightarrow{f^*(\beta_\xi)} \underbrace{f^*(\xi)}_{\hot_{S',W}(X')} \quad ,
\]
and check that the composite is the identity. We have set
\[
X' = X\times_S S',
\]
and the diagram is therefore also
\[
\begin{tikzcd}
X' \arrow[r, "{\alpha_{X'}}"] \arrow[rr, bend left, "{\id_{X'}}"] & \Phi_S(X')\times_S S' \arrow[r, "f^*(\beta_\xi)"] & X' % Reading verified against PDF leaf 50: the curved identity arrow is solid.
\end{tikzcd}
\]
in $(\Cat/S')(\Wd_{S'})^{-1}$. To check commutativity, we can compose with
\[
X\times_S S' = X' \xrightarrow{i_X\times_S\id_{S'}} \Phi_S(X')\times_S S' \qquad\qquad \text{(where } i_X : X\longrightarrow\Phi_S(X)\text{)}
\]
which is in $\Wd_{S'}$ (\footnote[71]{\textbf{NB} Since $X$ is smooth$/S$, and $\Phi_S(X)$ too, $X\longrightarrow\Phi_S(X)$ remains a colocal $W$-equivalence under any base change.}), now we check, from the general construction of $\beta_\xi=\beta_X$ (preceding page), that we have commutativity in
\[
\begin{tikzcd}
\Phi_S(X')\times_S S' \arrow[r, "f^*(\beta_\xi)"] \arrow[dr, "\Phi_S(\pr_1^X)\times_S S'"'] & X' \arrow[d, "i_X\times_S S'"] \\
 & \Phi_S(X)\times_S S'.
\end{tikzcd}
\]

\leavevmode\msp{page 145}We are reduced to the commutativity of
\[
\begin{tikzcd}
X\times_S S' = X' \arrow[r, "{\alpha_{X'}}"] \arrow[dr, "i_X\times_S S'"', "?"] & \Phi_S(X')\times_S S' \arrow[d, "\Phi_S(\pr_1^X)\times_S S'"] \\
 & \Phi_S(X)\times_S S'
\end{tikzcd}
\tag*{$(*)$}
\]
% foot: Version 9 novembre 2001

\end{document}
